% ================================================================
% SECTION 6 — CONCLUSION
% Target: ~1.5 pages
% ================================================================

\section{Conclusion}
\label{sec:conclusion}

This thesis set out to answer a deceptively simple question: by how
much did the April 2022 KwaZulu-Natal floods reduce freight volumes
on South Africa's NATCOR rail corridor, and for how long did the
effect persist?
Answering it rigorously required constructing a credible
counterfactual --- an estimate of what NATCOR volumes would have been
in the absence of the flood --- and demonstrating that this
counterfactual is robust to the choice of estimation method, the
composition of the donor pool, and the specification of the
pre-treatment window.

The main empirical finding is clear and consistent.
Across five structurally distinct SCM-family estimators applied to a
full five-corridor donor pool, the 2022 KZN floods caused a sustained
reduction in NATCOR monthly freight volume of approximately
$-0.157$\,MT per month, accumulating to a cumulative loss of
$-5.17$\,MT over the 33-month post-flood window.
Every estimator passes in-time placebo tests with zero exceedances,
and leave-one-out donor stability checks confirm that no single
corridor is driving the result.
The gap between observed and counterfactual NATCOR volume remains
negative across every post-flood month in the observation window,
providing no evidence of full recovery by December 2024 --- more than
two and a half years after the event.

A second, methodological finding is equally important.
When the analysis is restricted to a single donor corridor, the five
estimators produce results ranging from $-11.1$\,MT to $+3.7$\,MT
--- a spread of nearly 15\,MT that renders any single estimate
meaningless for policy purposes.
Once an adequate five-corridor donor pool is constructed, all five
methods converge to within 0.15\,MT of one another.
This empirical demonstration of donor pool inadequacy as a source
of estimator divergence is, to the author's knowledge, the first
of its kind in the transport economics literature, and provides a
concrete diagnostic for applied researchers: estimator disagreement
is a symptom of identification failure, not a reason to prefer one
algorithm over another.

\subsection*{Contributions Revisited}

The three contributions stated in Section~\ref{sec:contrib} have
been delivered.
Methodologically, this thesis provides a replicable five-method
benchmark for SCM-based freight impact analysis.
Empirically, it delivers the first causal estimate of the 2022 KZN
flood effect on NATCOR volumes using a counterfactual identification
strategy.
Diagnostically, it provides a worked example of how donor pool
inadequacy manifests in practice and how it can be detected and
corrected.

\subsection*{Directions for Future Research}

Three directions stand out as natural extensions of this work.

First, a commodity-level or origin--destination disaggregation of
the NATCOR volume series would allow researchers to determine whether
the flood effect is concentrated in specific freight types ---
containerised imports, agricultural bulk, or manufactured exports ---
and whether different commodities exhibit different recovery
trajectories.
Such analysis would require Transnet to make disaggregated operational
data available, which currently it does not.

Second, the donor pool could in principle be extended beyond South
African corridors to include freight rail data from comparable
emerging-market operators in East or Southern Africa.
A richer international donor pool would strengthen the parallel
trends assumption and could allow the SCM to be applied to smaller
corridor-level shocks that currently lack adequate domestic controls.

Third, complementing the SCM results with a dynamic treatment effect
model --- such as local projections or a structural time-series
model with an intervention component --- would characterise the
time profile of recovery more precisely than the static average
treatment effect reported here.
In particular, such models could test formally whether the gap
is narrowing, constant, or widening over time, and could provide
a probabilistic forecast of the date at which full volume recovery
might be expected.

\subsection*{Closing Remark}

South Africa's rail freight system operates at a strategic crossroads.
The 2022 KZN floods demonstrated that a single extreme weather event
can inflict a lasting, quantifiable reduction in corridor throughput
that persists well beyond the physical repair window.
As climate models project increasing frequency and intensity of
extreme precipitation events across southern Africa, the economic
case for pre-emptive infrastructure resilience investment on corridors
like NATCOR will only strengthen.
The methods developed in this thesis provide a rigorous, replicable
framework for quantifying such impacts --- one that can inform
cost-benefit analyses, insurance pricing, and policy design as
South Africa's freight rail system navigates the decades ahead.
